\documentclass[10pt,a4paper]{amsart}
\usepackage[margin=19mm]{geometry}
\usepackage{amsmath,amssymb,mathtools}
\usepackage[scr=rsfs,cal=euler]{mathalfa}
\usepackage{xfrac}
\usepackage[unicode,hidelinks,bookmarks=false]{hyperref}
\newcommand{\ie}{i.e.\ }
\newcommand{\cf}{cf.\ }
\newcommand{\resp}{respectively\ }
\newcommand{\Subsection}{\subsection}
\providecommand{\nobreakdash}{\nobreak\hbox{-}}
\setlength{\emergencystretch}{2em}
\multlinegap0pt
\usepackage{bm}
\usepackage{bbm}
\usepackage{dsfont}
\usepackage{xspace}
\usepackage{cancel}
\usepackage{mathtools}
\usepackage{amsthm}
\makeatletter
\@ifundefined{theorem}{\newtheorem{theorem}{Theorem}}{}
\makeatother
\title[Tensor completions of 2-nilpotent finitely generated torsion]{Tensor completions of 2-nilpotent finitely generated torsion-free groups}
\author{Mikheil Amaglobeli, Alexei Miasnikov}
\date{Source: arXiv:2601.17600v1 (CC BY 4.0)}
\begin{document}
\maketitle

\noindent\emph{Source and licence.} Tensor completions of 2-nilpotent finitely generated torsion-free groups. Version arXiv:2601.17600v1 (CC BY 4.0), distributed under the Creative Commons Attribution 4.0 International licence, as recorded in the arXiv licence metadata for this exact version and shown on the abstract page https://arxiv.org/abs/2601.17600v1. Full rights notice: licenses/arxiv-2601.17600-CC-BY-4.0.txt.\par\smallskip

\noindent\textit{Editorial scope.} Counted unit: the Theorem whose source label is \texttt{th:tensor-completions} in the article (source0.tex lines 870--877, source label \texttt{th:tensor-completions}), delivered together with the article's own complete original proof of it (source0.tex lines 879--938, one \texttt{proof} environment of 60 lines). No mathematics is written, repaired or re-derived: statement and proof are the article's own text line for line. Required context retained uncounted: none. Every definition, display and label of those lines is retained unchanged. Omitted from this package and not redistributed: 1--869, 878--878, 939--1444. Nothing inside a retained excerpt is omitted or altered: every retained body is delivered exactly as the article's lines are written, so no omission receipt of any kind accompanies this package. Citations: the retained excerpts carry no citation command at all, so no bibliography entry of the article is transcribed and no reference is redistributed. Reference adaptation: none was needed and none was made; every reference inside the delivered unit resolves inside it, so no unresolved reference or citation is produced. Printed slips kept literal: no printed word, symbol, hypothesis or number of the counted statement or of its proof was altered. Layout and class: the article's own class is \texttt{article} with packages including amsmath, amssymb, amsthm, eucal, graphics, inputenc, babel, amsfonts, amsbsy, latexsym, cite, csquotes, amscd, hyperref; the wrapper is the lane's pinned \texttt{amsart} editorial wrapper at 10pt on A4, which restates the theorem-like environments the retained text needs and adds the article's own macro definitions that the retained text needs (none), with extra wrapper packages (bm, bbm, dsfont, xspace, cancel, mathtools). The delivered numbering is the unit's own. Assets: no retained span contains a figure environment, a graphics-inclusion command, a PDF-page inclusion, a table or a TikZ picture; no image file is copied into this package, the package declares no asset dependency and no image file is redistributed with it. Licence: version arXiv:2601.17600v1 of this article is distributed under the Creative Commons Attribution 4.0 International licence, as recorded in the arXiv licence metadata for this exact version and shown on the abstract page https://arxiv.org/abs/2601.17600v1. A full rights notice is delivered at licenses/arxiv-2601.17600-CC-BY-4.0.txt. Characters: every retained excerpt is pure ASCII, so the delivered engine, XeLaTeX with its default Latin Modern text font, reports no missing character.
\begin{theorem} \label{th:tensor-completions}
	Let $G$ be a finitely generated 2-nilpotent torsion-free group, $R$ be a binomial domain of characteristic zero, and $\mu:G \otimes R \to G \otimes_{\mathcal{H}} R$ be the canonical $R$-epimorphism. Let $\bar u \cdot \bar v$ be the Mal'tsev basis of $G$, and $H_{\bar u \cdot \bar v}$ be a subgroup of $G \otimes R$, as above. Then the following holds:
	\begin{itemize}
		\item [1)] The restriction of $\mu$ to $H_{\bar u \cdot \bar v}$ is an isomorphism of abstract groups $H_{\bar u \cdot \bar v}$ onto $G \otimes_{\mathcal{H}} R$. In particular, $G \otimes_{\mathcal{H}} R$ embeds in $G \otimes R$ as an abstract group, and $\mu$ is a retraction to $G \otimes_{\mathcal{H}} R$, viewed as a subgroup $H_{\bar u \cdot \bar v}$ of $G \otimes R$.
		\item [2)] The kernel $D = \ker \mu$ of $\mu$ is an abelian $R$-subgroup of $G \otimes R$ generated as an $R$-group by all c-commutators $c(g,h)_\alpha, g,h \in G \otimes R$, $\alpha \in R$.
		\item [3)] $G \otimes R \simeq (G \otimes_{\mathcal{H}} R) \times D$ as a direct product of abstract groups.
	\end{itemize}
\end{theorem}

\begin{proof}
	Note that the $R$-epimorphism $\mu: G\otimes R \to G \otimes_{\mathcal{H}} R$ is the identity on $G$, hence it maps $\bar u^{\bar\alpha} \cdot \bar v^{\bar \beta}$ of $G \otimes R$ to $\bar u^{\bar\alpha} \cdot \bar v^{\bar \beta}$ of $G \otimes_{\mathcal{H}} R$. Consequently, the restriction of $\mu$ to $H_{\bar u \cdot \bar v}$ is an isomorphism of abstract groups $H_{\bar u \cdot \bar v} \to G \otimes_{\mathcal{H}} R$. This proves 1).
	
	
	Therefore, on the one hand
	$$
	\mu((xy)^\lambda) = (\mu(x)\mu(y))^\lambda= \mu(x)^\lambda \mu(y)^\lambda [\mu(x),\mu(y)]^{-\binom{\lambda}{2}}.
	$$
	On the other hand
	$$
	\mu((xy)^\lambda) = \mu(x^\lambda y^\lambda [x,y]^{-\binom{\lambda}{2}} c(x,y)_\lambda)= \mu(x)^\lambda \mu(y)^\lambda [\mu(x),\mu(y)]^{-\binom{\lambda}{2}} \mu(c(x,y)_\lambda).
	$$
	It follows that $\mu(c(x,y)_\lambda) = e$, so $c(x,y)_\lambda \in \ker \mu$ for any $x,y \in G\otimes R$ and $\lambda \in R$.
    Denote by $D$ the $R$-subgroup of $G \otimes R$ generated by all c-commutators $c(x,y)_\lambda$, i.e.
	$$
	D = \langle c(x,y)_\lambda \mid x,y \in G\otimes R, \lambda \in R\rangle_R.
	$$
	It turns out that $D$ is a central $R$-subgroup of $G \otimes R$ and $D \leq \ker \mu$.
    We claim that $D = \ker \mu$ and $G\otimes R = H_{\bar u \cdot \bar v} \times D$ as abstract groups. Since $\mu$ is a retract onto $H_{\bar u \cdot \bar v}$ and $D \leq Z(G\otimes R)$, it suffices to show that $G\otimes R = H_{\bar u \cdot \bar v}\cdot D$.
	
	Now note that the group $G \otimes R$ \ $R$-generated by the group $G$. Therefore, each element $g \in G \otimes R$ can be represented as
	$$
	g = w(g_1, \ldots,g_k),
	$$
	where $w(z_1, \ldots,z_k)$ is an $R$-word and $g_1, \ldots,g_k \in G$. We prove by induction on the number $e(w)$ of exponentiations in $w$ that $g = h\cdot d$, where $h \in H_{\bar u \cdot \bar v}$ and $d \in D$. If $e(w) = 0$, then $g \in G \leq H_{\bar u \cdot \bar v}$. Since $H_{\bar u \cdot \bar v}\cdot D$ is closed under multiplication, it suffices to consider the case when
	$$
	g = (f_1^{\gamma_1} \ldots f_k^{\gamma_k})^\lambda,
	$$
	where $f_1, \ldots, f_k\in G \otimes R$ and $\gamma_1, \ldots, \gamma_k, \lambda \in R$.
    By induction, $f_i^{\alpha_i} \in H_{\bar u \cdot \bar v} \cdot D$, so their product also belongs to $H_{\bar u \cdot \bar v} \cdot D$. Therefore, we can assume that
	$$
	g = (\bar u^{\bar\alpha} \bar v^{\bar \beta} d)^\lambda = (\bar u^{\bar\alpha})^\lambda \bar v^{\bar \beta \lambda} d^\lambda,
	$$
	where $d \in D$, and $\bar \beta \lambda$ is the scalar product of $\bar \beta$ and $ \lambda$. Therefore, it suffices to show that $(\bar u^{\bar\alpha})^\lambda = (u_1^{\alpha_1}\ldots u_m^{\alpha_m})^\lambda\in H_{\bar u \cdot \bar v}\cdot D$. We use induction on $m$. Denote $f = u_1^{\alpha_1}\ldots u_{m-1}^{\alpha_{m-1}}$. Then
	$$
	(u_1^{\alpha_1}\ldots u_m^{\alpha_m})^\lambda = f^\lambda (u_m^{\alpha_m})^\lambda [f,u_m^{\alpha_m}]^{-\binom{\lambda}{2}} c(f,u_m^{\alpha_m})_\lambda.
	$$
By induction, $f^\lambda \in H_{\bar u \cdot \bar v} \cdot D$. Obviously, $(u_m^{\alpha_m})^\lambda = u_m^{\alpha_m\lambda} \in H_{\bar u \cdot \bar v}$ and $ c(f,u_m^{\alpha_m})_\lambda \in D$. Note that
$$
[f,u_m^{\alpha_m}]^{-\binom{\lambda}{2}} = [f,u_m]^{-\binom{\lambda}{2}\alpha_m}
$$
and
$$
[f,u_m] = [u_1^{\alpha_1}\ldots u_{m-1}^{\alpha_{m-1}},u_m] = \prod_{i=1}^{m-1} [u_i,u_m]^{\alpha_i},
$$
from which it follows that
$$
[f,u_m]^{-\binom{\lambda}{2}\alpha_m} = \left(\prod_{i=1}^{m-1} [u_i,u_m]^{\alpha_i}\right)^{-\binom{\lambda}{2}\alpha_m} = \prod_{i=1}^{m-1} [u_i,u_m]^{-\alpha_i\binom{\lambda}{2}\alpha_m}.
$$
Since $[u_i,u_m] \in Z(G)$, then
$$
[u_i,u_m]^{-\alpha_i\binom{\lambda}{2}\alpha_m} = \bar v^{\bar \delta_{i,j}}
$$
for some tuple $\bar \delta_{i,j} \in R^n$. Therefore,
$$
\prod_{i=1}^{m-1} [u_i,u_m]^{-\alpha_i\binom{\lambda}{2}\alpha_m} = \prod_{i=1}^{m} \bar v^{\bar \delta_{i,j}} = \bar v^{\delta}
$$
for some $\delta \in R^n$. This proves that $[f,u_m^{\alpha_m}]^{-\binom{\lambda}{2}} \in H_{\bar u \cdot \bar v}$ and, consequently, that $(\bar u^{\bar\alpha})^\lambda = (u_1^{\alpha_1}\ldots u_m^{\alpha_m})^\lambda\in H_{\bar u \cdot \bar v}\cdot D$. We obtain that $g = (f_1^{\gamma_1} \ldots f_k^{\gamma_k})^\lambda \in H_{\bar u \cdot \bar v}\cdot D$, and hence by induction $g = w(g_1, \ldots,g_k) \in H_{\bar u \cdot \bar v}\cdot D$. Consequently, $G \otimes R =H_{\bar u \cdot \bar v}\cdot D$ and $D = \ker \mu$. Whence
$H_{\bar u \cdot \bar v} \cap D = \{e\}$ and $G \otimes R =H_{\bar u \cdot \bar v}\times D$. This proves the theorem.
\end{proof}

\end{document}
